
\subsection{Background}

A long-standing empirical regularity is that income distributions are approximately log-normal \citep{aitchison1957lognormal}. Multiplicative processes provide a rationale: incomes result from the interaction of factors such as skills, local labor markets and productivity \citep{neal2000theories}, and if they grow by proportional shocks that are independent of their level, as in Gibrat's law of proportionate effect \citep{gibrat1931inegalites}, the logarithm of income is a sum of independent shocks and tends to normality. The approximation works best within groups exposed to similar shocks. \citet{battistin2009consumption} show that, within cohorts, consumption, which tracks permanent income, is closer to log-normal than current income, whose transitory component departs from log-normality. Mixtures of log-normals fit heterogeneous populations well, each component describing a relatively homogeneous subgroup \citep{flachaire2007, lubrano2016}, and they improve substantially on a single log-normal when mapping income distributions across small areas \citep{gardini2022}.

The literature on neighborhood effects suggests why census tracts are natural components of such a mixture. Contextual influences and endogenous social interactions, such as shared information about jobs, peer effects and common exposure to local shocks, make income-generating processes more alike within neighborhoods than between them \citep{manski1993identification, durlauf1996neighborhoods, durlauf2004neighborhood}.

The log-normal is known to fit the tails poorly. Top incomes follow a power law rather than a log-normal decay \citep{pareto1896cours, atkinson2011top, gabaix2016power}, and the lower tail of Spanish tract distributions is heavier than a log-normal with the same median and Gini coefficient implies. I therefore use a more flexible family, the generalized beta distribution of the second kind (GB2) of \citet{mcdonald1984some}. It nests the log-normal as a limiting case, together with other distributions commonly fitted to incomes, and has a Pareto upper tail \citep{kleiber2003statistical}; and it can be estimated from the kind of summary statistics that statistical offices publish \citep{chotikapanich2007estimating}.

\subsection{Tract distributions and their aggregation}

Let $Y$ be the income per consumption unit of a person living in tract $j$. I assume that $Y$ follows a GB2 distribution, with cumulative distribution function
\begin{equation}\label{eq:gb2}
F_j(y) = I_{t_j(y)}(p_j, q_j), \qquad t_j(y) = \frac{(y/b_j)^{a_j}}{1 + (y/b_j)^{a_j}},
\end{equation}
where $I_t(p, q)$ is the regularized incomplete beta function, $b_j > 0$ is a scale parameter and $a_j, p_j, q_j > 0$ are shape parameters. The GB2 nests the Singh--Maddala ($p_j = 1$) and Dagum ($q_j = 1$) distributions and, as a limiting case, the log-normal. Its upper tail is Pareto with index $a_j q_j$, which I constrain to exceed one so that the mean is finite, and I parameterize it by $\theta_j = (\ln a_j, \ln m_j, \ln p_j, \ln(a_j q_j - 1))$, where $m_j$ is the median.

I estimate $\theta_j$ separately for each tract by weighted minimum distance to the thirteen indicators that the ADRH publishes for it,
\begin{equation}\label{eq:md}
\hat\theta_j = \arg\min_{\theta} \sum_{k=1}^{13} \rho_{jk}\big(g_k(\theta) - g_{jk}\big),
\end{equation}
where $g_{jk}$ is the published value of indicator $k$ and $g_k(\theta)$ its value under the GB2: the median, the mean, the Gini coefficient (computed, as in the ADRH, after replacing the lowest and highest 1.5\% of incomes), the P80/P20 ratio and the nine shares. Differences are measured in log points for the median, the mean and the P80/P20 ratio, in points for the Gini coefficient and in percentage points for the shares, and $\rho_{jk}(d) = (d/s_{jk})^2$, where $s_{jk}$ combines the rounding of the published figure with a tolerance of \TolShare{} points for the shares, \TolGini{} points for the Gini coefficient, \TolMeanPct{}\% for the mean and \TolPPct{}\% for the P80/P20 ratio. Two features of the publication enter $\rho_{jk}$. Medians, published as the midpoints of \texteuro 700 intervals, are treated as intervals: inside the interval the pull towards the published value is weak ($s = \MedPullPct{}\%$) and outside it steep ($s = \MedOutPct{}\%$). Values at the ADRH's caps (Section \ref{sec:data}) are treated as censored: beyond the cap, the tolerance widens to \CapAllowPct{}\% of the capped value. The minimization uses a Levenberg--Marquardt algorithm started from Dagum and Singh--Maddala fits, themselves started from the log-normal with the tract's median and Gini coefficient; \GbConvergedPct{}\% of the \GbTracts{} tracts converge.

The shares are not published for \GbPriorTracts{} of these tracts. For them, the share targets are predictions from regressions of the logit of each published share on a cubic polynomial in the tract's log median, log ratio of mean to median, Gini coefficient and log P80/P20 ratio, and on log population, with province effects shrunk towards zero, estimated on the \GbPriorTrain{} tracts with published shares and no capped indicator. Their tolerance adds the regressions' root mean square error (\GbPriorRMSE{} percentage points on average). In the \GbGapMuns{} municipalities that publish shares, the predictions are shifted so that the municipality's tracts together reproduce its published shares.

The \LocMeanTracts{} tracts without these indicators (Section \ref{sec:data}) get a log-normal, $\ln Y \sim N(\nu_j, \sigma_j^2)$. Its Gini coefficient is $G_j = 2\,\Phi(\sigma_j/\sqrt{2}) - 1$, where $\Phi$ is the standard normal cumulative distribution function, so, given the tract's imputed Gini coefficient $G_j$ (Section \ref{sec:gini}) and imputed mean $\bar y_j$, I set
\begin{equation}\label{eq:params}
\sigma_j = \sqrt{2}\,\Phi^{-1}\!\left(\frac{G_j + 1}{2}\right), \qquad \nu_j = \ln \bar y_j - \frac{\sigma_j^2}{2}.
\end{equation}

The income distribution of an area $A$ (Spain, a province or a municipality) is the population-weighted mixture of the distributions of its tracts, with cumulative distribution function
\begin{equation}\label{eq:mixture}
F_A(y) = \sum_{j \in A} w_j\, F_j(y), \qquad w_j = \frac{N_j}{\sum_{k \in A} N_k},
\end{equation}
where $N_j$ is the population of tract $j$. The $p$-th percentile $q_{A,p}$, $p = 1, \dots, 99$, solves $F_A(q_{A,p}) = p/100$. I solve this equation with Brent's method, bracketing the root on a grid of incomes; the solutions satisfy it to within $10^{\RootErrExp}$. Of the \MunLookup{} municipalities covered, \MunSingleTract{} (\MunSingleTractPct{}\%, home to \MunSingleTractPopPct{}\% of the population) consist of a single tract, so their distribution is that tract's. The \MunDropped{} municipalities in which no tract has income data are not covered.

\subsection{Missing Gini coefficients}\label{sec:gini}

For the \GiniImputedTracts{} tracts with income data but no published Gini coefficient, which get the log-normal, I predict it with gradient-boosted regression trees (XGBoost) trained on the \GiniTrainTracts{} tracts with a published Gini. The predictors are the logarithm of mean income per consumption unit, the dependency ratio, mean age, the share of the population under 18, the share of single-person households, mean household size, population and province dummies. The model uses \XgbRounds{} boosting rounds, a maximum tree depth of \XgbDepth{} and a learning rate of \XgbEta{}, without tuning. Table \ref{tab:gini} compares it with OLS on the same predictors, with province fixed effects, and with a constant, using the same \CvFolds{}-fold cross-validation folds for all three. XGBoost reduces the out-of-sample root mean square error (RMSE) by \XgbRMSEGain{}\% relative to OLS, with an out-of-sample $R^2$ of \XgbRsq{}.

\begin{table}[H]
\centering
\begin{threeparttable}
\onehalfspacing
\caption{Gini imputation: out-of-sample fit}\label{tab:gini}
\footnotesize
\begin{tabular}{@{}lcccc@{}}
\toprule
Model & RMSE & MAE & MAPE (\%) & $R^2$ \\
\midrule
Constant (training mean) & 3.87 & 2.96 & 10.33 & 0.00 \\
OLS, province fixed effects & 2.94 & 2.25 & 7.87 & 0.42 \\
XGBoost, province dummies & 2.57 & 1.97 & 6.89 & 0.56 \\
\bottomrule
\end{tabular}

\begin{tablenotes}[para,flushleft]
\scriptsize
\textbf{Notes}: \CvFolds{}-fold cross-validation over the \GiniTrainTracts{} tracts with a published Gini coefficient, with the same folds for every model. RMSE is the root mean square error, MAE the mean absolute error and MAPE the mean absolute percentage error, all out of sample; $R^2$ is one minus the ratio of the out-of-sample squared error to the variance of the Gini coefficient. The Gini coefficient is measured in points (0--100).
\end{tablenotes}
\end{threeparttable}
\end{table}

Cross-validation measures accuracy for tracts like those in the training sample, all of which have at least \ObsGiniMinPop{} residents, whereas the tracts to impute are smaller. The model extrapolates poorly in that direction: trained on tracts with at least \SmallTrainMin{} residents, it predicts the Gini coefficient of the \SmallTestTracts{} tracts with \ObsGiniMinPop{} to \SmallTestMaxPop{} residents with an RMSE of \XgbSmallRMSE{}, no better than the standard deviation of the Gini across tracts (\GiniSD{}). The imputed Gini coefficients average \GiniPredMean{}, against \GiniObsMean{} for published ones. The imputation hardly affects the national and provincial distributions, since these tracts house \GiniMissPopPct{}\% of the population, but it matters for the smallest municipalities. In \MunSynthetic{} municipalities (\MunSyntheticPct{}\% of those covered, with \MunSyntheticPopPct{}\% of the population), no tract has a published Gini coefficient, so every tract gets the log-normal. Their income per person is, moreover, an average over small municipalities of the same agricultural district, which takes only \MunSyntheticIncomeValues{} distinct values across them. Their distributions are therefore model-based approximations rather than estimates specific to each municipality.

\subsection{Nowcasting}\label{sec:nowcast}

The ADRH is published almost two years after the income year, so its latest edition refers to \BaseYear{}. To compare households with a distribution of the same year, I project incomes to \TargetYear{} with the ECV, whose wave $t+1$ reports mean income per consumption unit in year $t$.

Let $g^{\text{ADRH}}_t$ and $g^{\text{ECV}}_t$ be the growth of national mean income per consumption unit in year $t$ in each source, the ADRH mean being the population-weighted average over the municipalities observed in both years. Survey and administrative growth rates need not coincide, so I measure how ECV growth has compared with ADRH growth over the years both cover,
\begin{equation}
\rho = \frac{\displaystyle\sum_{t=\NowcastFirstYear}^{\NowcastLastYear} g^{\text{ADRH}}_t}{\displaystyle\sum_{t=\NowcastFirstYear}^{\NowcastLastYear} g^{\text{ECV}}_t},
\end{equation}
and nowcast growth in \TargetYear{} as
\begin{equation}
\hat{g}_{\TargetYear} = \rho \, g^{\text{ECV}}_{\TargetYear}.
\end{equation}
Over \NowcastFirstYear--\NowcastLastYear{} the ECV slightly overstates administrative income growth ($\rho = \NowcastRho{}$). With ECV growth of \EcvGrowth{}\% in \TargetYear{}, the nowcast growth is \NowcastGrowth{}\%.

I scale the income of every tract by $1+\hat{g}_{\TargetYear}$. Multiplying a GB2 variable by a constant multiplies its scale $b_j$ by that constant and leaves $a_j$, $p_j$ and $q_j$ unchanged (for a log-normal, it shifts $\nu_j$ by $\ln(1+\hat{g}_{\TargetYear})$ and leaves $\sigma_j$ unchanged), so the adjustment moves every distribution while preserving its shape and Gini coefficient. The underlying assumption is that incomes grew at a common rate across areas and along the distribution over one year.

The projection is uncertain. In leave-one-year-out validation, $\rho\, g^{\text{ECV}}_t$ misses $g^{\text{ADRH}}_t$ by \NowcastLooMAE{} percentage points on average, against \EcvRawMAE{} points for uncorrected ECV growth. The estimate of $\rho$ also depends on the window and the estimator: it is \RhoLastThree{} over the last three years, \RhoExclCovid{} excluding 2020 and 2021, and \RhoOrigin{} by least squares through the origin. Across these choices, including no correction, nowcast growth ranges from \NowcastGrowthMin{}\% to \NowcastGrowthMax{}\%.

\subsection{Placing a household}\label{sec:placing}

\href{https://comparatuingreso.es/}{\mbox{Comparatuingreso.es}} asks for the household's net monthly income in \TargetYear{}, the year of the nowcast distribution, and for whether it was paid in 12 or 14 instalments (Spanish salaries and pensions are often paid in 14, with two extra payments a year). Annual income is the monthly amount times 12 or 14. Dividing it by the household's consumption units on the modified OECD scale, from the number of members aged 14 or over and under 14, gives income per consumption unit, $y$.

The household's position in area $A$ is
\begin{equation}
P_A(y) =
\begin{cases}
0 & \text{if } y < q_{A,1}, \\
\max\{p \in \{1, \dots, 99\} : q_{A,p} \le y\} & \text{otherwise.}
\end{cases}
\end{equation}
Since $F_A(y) \ge P_A(y)/100$, a household with $P_A(y) \ge 1$ has more income per consumption unit than at least $P_A(y)\%$ of the population, which is how the app reports it; $P_A(y) = 0$ is reported as being among the 1\% with the lowest income. Positions are shares of persons, not households: every person counts with the income per consumption unit of their household.
